\documentclass[10pt]{article}
\usepackage[margin=1in]{geometry}
\usepackage{booktabs}
\usepackage{array}
\usepackage{multirow}
\usepackage{amsmath}
\usepackage{graphicx}
\usepackage{hyperref}
\usepackage{parskip}
\usepackage{titlesec}
\usepackage{caption}
\usepackage{float}

% Tighten spacing
\titlespacing*{\section}{0pt}{6pt}{3pt}
\titlespacing*{\subsection}{0pt}{4pt}{2pt}
\setlength{\parskip}{3pt}
\setlength{\parindent}{0pt}
\captionsetup{font=small, skip=3pt}

\begin{document}

%% ----------------------------------------------------------------
%% TITLE PAGE
%% ----------------------------------------------------------------
\begin{titlepage}
\centering
\vspace*{2cm}
{\LARGE \textbf{MAE 155B --- Aircraft Design}} \\[0.4cm]
{\Large \textbf{Detail Design Review Report}} \\[0.8cm]
\rule{\linewidth}{0.5pt} \\[0.4cm]
{\large \textbf{Group 2 --- Package Delivery UAV}} \\[0.4cm]
\rule{\linewidth}{0.5pt} \\[0.8cm]

\begin{tabular}{ll}
\textbf{Team Members:} & Harshil Patel \\
                       & John Sigafoos \\
                       & Angel Ochoa \\
                       & Juan Sanchez \\
                       & Sara Chowdhury \\
                       & Analisa Veloz \\[0.5cm]
\textbf{Course:}       & MAE 155B --- Aircraft Design \\
\textbf{Date:}         & May 12, 2026 \\
\textbf{Report Type:}  & Detail Design Review (DDR) \\
\end{tabular}

\vspace{1.5cm}
\textbf{Abstract} \\[0.2cm]
\begin{minipage}{0.85\linewidth}
\normalsize
This report presents the Detail Design of a small fixed-wing RC aircraft for autonomous
package delivery, optimized to maximize profit per unit time. The final design is a
tailless flying wing with twin delta winglets, an 8.0 AR tapered planform, and an
APC 10$\times$4.5 MR propeller driven by a 1100\,KV brushless motor. The
design achieves a projected profit rate of \$11.04/hr at an 800\,g payload,
a static margin of 11.05\%\,MAC, and a total mission energy consumption of
16.5\,Wh within the 18.98\,Wh available (28\% reserve).
\end{minipage}
\end{titlepage}

%% ----------------------------------------------------------------
%% EXECUTIVE SUMMARY
%% ----------------------------------------------------------------
\section*{Executive Summary}

The Group 2 aircraft is a flying-wing UAV designed for short-range package delivery
on a closed course. Design optimization maximizes profit per unit time $J$ by
balancing payload capacity, lift-to-drag ratio, and energy consumption against
structural and stability constraints.

\subsection*{Top-Level Specifications}

\begin{table}[H]
\centering
\caption{Top-Level Design Specifications (DDR)}
\label{tab:specs}
\small
\begin{tabular}{lll}
\toprule
\textbf{Parameter} & \textbf{Value} & \textbf{Unit} \\
\midrule
\multicolumn{3}{l}{\textit{Geometry}} \\
Wingspan (full) & 1.5715 & m \\
Wing area & 0.3087 & m$^2$ \\
Aspect ratio & 8.0 & --- \\
Taper ratio & 0.661 & --- \\
Quarter-chord sweep & 28.3 & deg \\
Root / Tip chord & 0.2365 / 0.1563 & m \\
MAC & 0.1992 & m \\
\midrule
\multicolumn{3}{l}{\textit{Mass \& Weights}} \\
MTOW (loaded) & 2.750 & kg (27.0\,N) \\
Empty weight (no payload) & 1.950 & kg \\
Payload & 800 & g \\
\midrule
\multicolumn{3}{l}{\textit{Performance}} \\
Design cruise speed & 20.0 & m/s \\
Stall speed (loaded) & 12.7 & m/s \\
Takeoff speed & 14.2 & m/s \\
Wing loading ($W/S$) & 87.4 & N/m$^2$ \\
Static thrust & 15.6 & N \\
Static thrust-to-weight & 0.58 & --- \\
Min turn radius & 32 & m \\
Max bank angle & 52 & deg \\
\midrule
\multicolumn{3}{l}{\textit{Propulsion}} \\
Motor & 1100\,KV brushless & --- \\
Propeller & APC 10$\times$4.5 MR & --- \\
Battery & 3S LiPo 11.1\,V & --- \\
Mission energy & 16.5 & Wh \\
Energy with reserve & 19.0 & Wh \\
\midrule
\multicolumn{3}{l}{\textit{Stability}} \\
Static margin (AVL) & 11.05 & \% MAC \\
CG location (loaded) & 16.7 & \% MAC \\
Trim elevon & 3.68 & deg \\
Short-period $\omega_n$ / $\zeta$ & 9.0 / 0.46 & rad/s / --- \\
Dutch roll $\omega_n$ / $\zeta$ & 5.3 / 0.09 & rad/s / --- \\
\midrule
\multicolumn{3}{l}{\textit{Economics}} \\
Profit rate (800\,g payload) & \textbf{11.04} & \$/hr \\
\bottomrule
\end{tabular}
\end{table}

\noindent\textit{[Technical 3-view drawing to be inserted here from SolidWorks/CAD export.
Key dimensions: span = 1.572\,m, fuselage length = 0.950\,m,
fin height = 0.172\,m at $y = 0.860$\,m from CL.]}

%% ----------------------------------------------------------------
%% BODY
%% ----------------------------------------------------------------
\section{Conceptual Design Phase}

\subsection{Mission Requirements and Scoring Function}

The design objective is to maximize profit per unit time $J$ (\$/hr), defined as:

\begin{equation}
J = \frac{P_r \cdot V_p \cdot V_{ps}}{T_f \cdot V_{p,\text{ref}}}
\end{equation}

where $P_r$ is the per-delivery revenue, $V_p = 0.015$\,m$^3$ is the payload volume,
$V_{ps} = 15$ is the volume scaling ratio, and $T_f$ is total mission time.
The mission profile consists of runway takeoff ($\sim$140\,m), a climb to altitude,
cruise laps on a 509\,m circuit, and a landing recovery.

\subsection{Configuration Selection}

A tailless flying wing was selected to minimize wetted area and parasite drag,
maximizing $L/D$ for a given wing loading. The configuration eliminates
horizontal tail download, improving payload fraction relative to conventional
tractor designs at small scale. A twin-boom or conventional aft-tail was
rejected because the added wetted area reduces cruise $L/D$ below the
$\approx$14 required for the target profit rate.

Vertical stability is provided by twin delta winglets mounted at the wing tips,
which contribute directional stability ($C_{n\beta} > 0$) and house split rudder
surfaces for active yaw control.

\section{Preliminary Design Phase}

\subsection{Sizing and Constraint Analysis}

Wing loading was sized from the stall constraint with $C_{L,\text{max}} = 0.915$
and a target stall speed of 12.0\,m/s, yielding $W/S \leq 79.6$\,N/m$^2$.
Thrust loading was sized from the takeoff constraint (governing), giving
$T/W \geq 0.47$.

A parametric sweep of span $b$, taper $\lambda$, and payload mass $m_p$ was
performed through the full MATLAB analysis pipeline with AVL in the loop.
The profit-optimal solution at the static-margin constraint $\text{SM} \geq 5\%$
is a 1200\,g payload at $J = 12.83$\,\$/hr; however, the team selected
$m_p = 800$\,g ($J = 11.04$\,\$/hr) to provide SM = 9.2\% with less structural
loading risk during the prototype test campaign.

\subsection{Airfoil Selection}

The MH95 airfoil was selected for the centerbody cross-section due to its
favorable reflexed camber line, which naturally provides the positive pitching
moment coefficient ($C_{m0} > 0$) required for pitch trim in a tailless
configuration without large positive elevon deflections.
Wing panels use E222 at the root and E230 at the tip,
both generating low $C_{D,\text{profile}}$ at Reynolds numbers of
$2.1 \times 10^5$--$3.1 \times 10^5$.
The 3D wing $C_{L,\text{max}} = 0.915$ was verified via lifting-line analysis.

\subsection{Propulsion Sizing}

Motor and propeller selection was driven by the static thrust requirement
($T_{\text{static}} \geq T/W \cdot W_0 \approx 12.7$\,N).
An APC 10$\times$4.5 MR propeller on a 1100\,KV motor at 11.1\,V produces
$T_{\text{static}} = 15.6$\,N at $I_{\text{static}} = 26.7$\,A,
below the 35\,A ESC limit.
The motor--propeller torque balance was solved numerically using APC tabulated
performance data across the full flight speed envelope.

\section{Detail Design Phase}

\subsection{Aerodynamics}

Panel aerodynamics and lift distribution were computed using AVL (Athena Vortex
Lattice). Key cruise aerodynamic coefficients:

\begin{table}[H]
\centering
\caption{Cruise Aerodynamic Performance at $V = 20$\,m/s}
\small
\begin{tabular}{lll}
\toprule
Parameter & Value & \\
\midrule
$C_{D,0}$ (total) & 0.0183 & (wing 0.0127, fuse 0.0040, fins 0.0016) \\
$C_L$ (cruise) & 0.350 & \\
$\alpha$ (cruise) & 2.46 & deg \\
$L/D$ (cruise) & 14.34 & \\
$(L/D)_{\text{max}}$ & 16.57 & \\
$C_{L,\alpha}$ (3D) & 0.090 & /deg \\
\bottomrule
\end{tabular}
\end{table}

The elevon chord fraction of 45\% and span $\eta = 0.60$--$0.95$ was set by the
hinge-moment constraint: the SG90 servo (0.177\,N$\cdot$m stall torque) is loaded
to 40\% capacity at maximum 20$^\circ$ deflection.

\subsection{Structural Design}

The primary structural member is a 10\,mm-diameter carbon fiber spar tube
that carries the full bending moment at the root. At the 2.0$g$ design load
with a safety factor of 2.0, the required spar diameter is 6.36\,mm;
the 10\,mm tube provides a bending factor of safety of 182 and a predicted
wingtip deflection of 0.5\,mm (0.03\% semispan). The wing skin uses 3\,mm
foam composite sheeting. The fuselage is a 0.95\,m monolithic foam-core
shell shaped to the MH95 cross-section; its CAD-estimated mass is 1135\,g.

\subsection{Stability and Control}

\textbf{Static Stability.}
The neutral point was computed by AVL. The CG is located at 16.7\%\,MAC (loaded),
giving SM = 11.05\%\,MAC, satisfying the $\geq 5\%$ requirement with margin.

\textbf{Longitudinal Dynamic Modes.}
\begin{itemize}\itemsep0pt
  \item Short period: $\omega_n = 9.0$\,rad/s, $\zeta = 0.46$ — well-damped, satisfactory.
  \item Phugoid: $\omega_n = 0.42$\,rad/s, $\zeta = -0.031$ — slightly unstable;
    time to double amplitude $\approx 25$\,s, easily controlled by pilot-in-the-loop.
\end{itemize}

\textbf{Lateral-Directional Modes.}
\begin{itemize}\itemsep0pt
  \item Dutch roll: $\omega_n = 5.3$\,rad/s, $\zeta = 0.091$ — marginally damped.
    The twin winglet fins provide $C_{n\beta} = 0.038$/rad (positive directional stability).
    Active rudder deflection up to $\delta_r = 25^\circ$ trims up to $\beta = 8.6^\circ$
    of sideslip.
  \item Aileron roll rate at $\delta_a = 20^\circ$: 238\,deg/s — aggressive authority.
\end{itemize}

\textbf{Control Surfaces.}
Elevons (pitch + roll) and twin rudders (yaw) are each driven by SG90 servos
confirmed within torque limits:
elevon HM = 0.070\,N$\cdot$m (40\% of capacity),
rudder HM = 0.022\,N$\cdot$m (13\% of capacity).
Trim elevon deflection at cruise is $\delta_e = +3.68^\circ$ (trailing-edge down).

\subsection{Energy and Mission Analysis}

\begin{table}[H]
\centering
\caption{Mission Energy Budget}
\small
\begin{tabular}{lll}
\toprule
Segment & Energy & Notes \\
\midrule
Climb & 4113\,J (1.14\,Wh) & $V_{\text{climb}} = 11$\,m/s, $\gamma = 6^\circ$ \\
Cruise & 55293\,J (15.36\,Wh) & 1 lap, $V = 20$\,m/s \\
Total mission & 59406\,J (16.50\,Wh) & \\
Design (with 15\% reserve) & 68316\,J (18.98\,Wh) & 28\% margin above mission \\
\bottomrule
\end{tabular}
\end{table}

\subsection{Economics}

The profit function uses the course-defined scoring:
$J = 11.04$\,\$/hr at 800\,g payload, $W_g = 2750$\,g.
The unconstrained optimum is \$12.83/hr at 1200\,g payload; the 800\,g selection
trades \$1.79/hr for 2.9\% additional static margin, justified by prototype safety
during the initial test campaign.

\section{Component Fabrication Plan}

\begin{table}[H]
\centering
\caption{Component Fabrication Summary}
\small
\begin{tabular}{lll}
\toprule
Component & Method & Material \\
\midrule
Fuselage/centerbody & CNC hot-wire + foam carving & EPS foam, fiberglass skin \\
Wing panels & Foam sheeting + rib alignment & 3\,mm XPS foam, CF tape LE \\
Main spar & COTS tube & 10\,mm OD carbon fiber \\
Winglet fins & Foam cut to profile & 3\,mm XPS foam \\
Servos (5$\times$) & COTS & SG90, 9\,g \\
Motor mount & 3D printed & PLA/PETG \\
Cargo bay & Foam lined cavity & EPS w/ foam lid \\
Landing gear & Bent wire + foam skid & Music wire + EPS \\
\bottomrule
\end{tabular}
\end{table}

Key fabrication sequences: (1) cut and assemble foam fuselage shell;
(2) insert 10\,mm spar through fuselage spar tunnel;
(3) attach wing panels and verify CG at $x = 0.337$\,m from motor;
(4) install servos, receiver, and ESC;
(5) balance battery to $x = 0.61$\,m to achieve $\text{SM} = 11\%$.
A static weight-and-balance check before first flight will confirm CG
within 14--20\%\,MAC.

%% ----------------------------------------------------------------
%% REFERENCES
%% ----------------------------------------------------------------
\section*{References}

\begin{enumerate}\itemsep0pt
  \item Drela, M. and Youngren, H., \textit{AVL 3.36 User Primer}, MIT, 2017.
  \item Selig, M.S., \textit{UIUC Airfoil Database}, University of Illinois, 1995.
  \item APC Propellers, \textit{Performance Data Files --- PER3\_10x45MR}, Landing Products, 2023.
  \item Raymer, D.P., \textit{Aircraft Design: A Conceptual Approach}, 6th ed., AIAA, 2018.
  \item MAE 155B Course Slides, UCSD, Spring 2026 (Sizing, Stability, Propulsion).
\end{enumerate}

%% ----------------------------------------------------------------
%% APPENDIX
%% ----------------------------------------------------------------
\appendix
\section{Profit Optimization Sensitivity}

\begin{table}[H]
\centering
\caption{Profit vs. Payload at Constant Geometry ($b = 1.572$\,m, $AR = 8$)}
\small
\begin{tabular}{cccc}
\toprule
Payload (g) & $J$ (\$/hr) & SM (\% MAC) & Note \\
\midrule
600 & 8.92 & 12.1 & Under-payload \\
800 & 11.04 & 9.2 & \textbf{Selected} \\
1000 & 11.87 & 7.1 & Acceptable SM \\
1200 & 12.83 & 5.0 & Optimizer optimum \\
1400 & 13.31 & 2.4 & SM constraint violated \\
\bottomrule
\end{tabular}
\end{table}

\section{Control Derivative Summary (AVL)}

\begin{table}[H]
\centering
\caption{AVL Stability and Control Derivatives at Cruise}
\small
\begin{tabular}{llll}
\toprule
Derivative & Value & Units & Notes \\
\midrule
$C_{L\alpha}$ & 0.0902 & /deg & 3D lifting line \\
$C_{m\alpha}$ & $-0.338$ & /rad & Pitch stiffness (stable) \\
$C_{n\beta}$ & 0.038 & /rad & Directional stability (stable) \\
$C_{l\delta_a}$ & 0.0062 & /deg & Aileron roll effectiveness \\
$C_{n\delta_r}$ & 0.000287 & /deg & Rudder yaw effectiveness \\
$C_{l p}$ & $-0.496$ & /rad & Roll damping \\
$C_{m q}$ & $-4.06$ & /rad & Pitch damping \\
\bottomrule
\end{tabular}
\end{table}

\end{document}
